\documentclass[border=12pt]{standalone}
\usepackage[no-math]{fontspec}
\usepackage[brazil]{babel}
\usepackage{amsmath,amssymb}
\IfFileExists{C:/Windows/Fonts/Lato_Regular.ttf}{%
  \setmainfont{Lato}[
    Path = C:/Windows/Fonts/,
    UprightFont = Lato_Regular.ttf,
    BoldFont = Lato_Bold.ttf,
    ItalicFont = Lato_Italic.ttf,
    BoldItalicFont = Lato_BoldItalic.ttf
  ]
  \setsansfont{Lato}[
    Path = C:/Windows/Fonts/,
    UprightFont = Lato_Regular.ttf,
    BoldFont = Lato_Bold.ttf,
    ItalicFont = Lato_Italic.ttf,
    BoldItalicFont = Lato_BoldItalic.ttf
  ]
}{%
  \RequirePackage[default]{lato}
}
\renewcommand{\familydefault}{\sfdefault}

\usepackage{xcolor}
\usepackage{tikz}
\usepackage{pgfplots}
\pgfplotsset{compat=1.18}
\usetikzlibrary{arrows.meta,positioning,shapes.geometric,fit,backgrounds,calc,decorations.pathreplacing,patterns}

% Paleta institucional UEFS / Extensão Rural
\definecolor{erNavy}{HTML}{1F4268}
\definecolor{erNavyDark}{HTML}{152E4A}
\definecolor{erGold}{HTML}{B38A3D}
\definecolor{erGoldDark}{HTML}{8C6927}
\definecolor{erTeal}{HTML}{0E7C7B}
\definecolor{erForest}{HTML}{2D6A4F}
\definecolor{erAlert}{HTML}{9E2A2B}
\definecolor{erInk}{HTML}{26313D}
\definecolor{erMuted}{HTML}{5B6573}

\definecolor{bgNavy}{HTML}{EDF2F7}
\definecolor{bgGold}{HTML}{F8F3E9}
\definecolor{bgTeal}{HTML}{EBF6F5}
\definecolor{bgForest}{HTML}{EEF5F0}
\definecolor{bgAlert}{HTML}{FDF2F2}

\tikzset{
  >=Stealth,
  erBox/.style={
    draw=#1, fill=white, rounded corners=5pt, line width=1.15pt,
    inner sep=7pt, align=center, font=\footnotesize\sffamily, text=erInk
  },
  erTint/.style={
    draw=#1, fill=#1!8!white, rounded corners=5pt, line width=1.15pt,
    inner sep=7pt, align=center, font=\footnotesize\sffamily, text=erInk
  },
  erBadge/.style={
    fill=#1, text=white, rounded corners=9pt, inner xsep=8pt, inner ysep=3pt,
    font=\scriptsize\bfseries\sffamily
  },
  erTitle/.style={font=\normalsize\bfseries\sffamily, text=erNavy, align=center},
  erSub/.style={font=\footnotesize\sffamily, text=erMuted, align=center},
  erArrow/.style={->, >=Stealth, line width=1.35pt},
  erDArrow/.style={<->, >=Stealth, line width=1.35pt}
}

\hyphenpenalty=10000
\exhyphenpenalty=10000

\begin{document}
\begin{tikzpicture}[font=\sffamily]

% Título principal
\node[erTitle, font=\large\bfseries\sffamily] (title) at (0, 5.05)
  {Bases Pedagógicas da Extensão Rural: Knowles, Freire e Kolb};

% Card 1: Andragogia (Malcolm Knowles)
\node[
  draw=erNavy, fill=bgNavy, rounded corners=6pt, line width=1.25pt,
  text width=4.45cm, inner sep=10pt, align=left, anchor=north
] (k1) at (-5.6, 4.35) {%
  \parbox[t][3.6cm][t]{4.45cm}{%
    {\centering{\color{erNavy}\small\bfseries 1. Andragogia\\(Malcolm Knowles)}\par\vspace{5pt}}
    {\color{erNavy!40}\rule{4.45cm}{0.8pt}}\par\vspace{8pt}
    {\raggedright\color{erInk}\footnotesize
    \textbf{\color{erNavy}\textbullet}\enspace Adultos aprendem pela experiência\\[7pt]
    \textbf{\color{erNavy}\textbullet}\enspace Autonomia e aplicação imediata ao problema rural}
  }%
};

% Card 2: Pedagogia Libertadora (Paulo Freire)
\node[
  draw=erTeal, fill=bgTeal, rounded corners=6pt, line width=1.25pt,
  text width=4.45cm, inner sep=10pt, align=left, anchor=north
] (k2) at (0, 4.35) {%
  \parbox[t][3.6cm][t]{4.45cm}{%
    {\centering{\color{erTeal}\small\bfseries 2. Pedagogia Libertadora\\(Paulo Freire)}\par\vspace{5pt}}
    {\color{erTeal!40}\rule{4.45cm}{0.8pt}}\par\vspace{8pt}
    {\raggedright\color{erInk}\footnotesize
    \textbf{\color{erTeal}\textbullet}\enspace Educação problematizadora (círculos de cultura)\\[7pt]
    \textbf{\color{erTeal}\textbullet}\enspace Superação da educação bancária}
  }%
};

% Card 3: Aprendizagem Experiencial (David Kolb)
\node[
  draw=erForest, fill=bgForest, rounded corners=6pt, line width=1.25pt,
  text width=4.45cm, inner sep=10pt, align=left, anchor=north
] (k3) at (5.6, 4.35) {%
  \parbox[t][3.6cm][t]{4.45cm}{%
    {\centering{\color{erForest}\small\bfseries 3. Aprendizagem Experiencial\\(David Kolb)}\par\vspace{5pt}}
    {\color{erForest!40}\rule{4.45cm}{0.8pt}}\par\vspace{8pt}
    {\raggedright\color{erInk}\footnotesize
    \textbf{\color{erForest}Ciclo formativo:}\\[5pt]
    Experiência Concreta $\rightarrow$\\[2pt]
    Observação Reflexiva $\rightarrow$\\[2pt]
    Conceituação $\rightarrow$\\[2pt]
    Experimentação Ativa}
  }%
};

% Base unificadora inferior
\node[
  draw=erGoldDark, fill=bgGold, rounded corners=6pt, line width=1.35pt,
  text width=15.65cm, minimum height=0.95cm, inner sep=9pt, align=center,
  font=\small\itshape\sffamily, text=erInk
] (base) at (0, -1.15) {%
  \textbf{\upshape\color{erGoldDark}Síntese formativa para a ATER:} articular a prática da lavoura com reflexão crítica coletiva.
};

% Conectores dos 3 pilares à base
\draw[erArrow, color=erNavy!75]   (k1.south) -- (k1.south |- base.north);
\draw[erArrow, color=erTeal!75]   (k2.south) -- (k2.south |- base.north);
\draw[erArrow, color=erForest!75] (k3.south) -- (k3.south |- base.north);

\end{tikzpicture}
\end{document}
